%================================================
% RAPPORT PHASE 3 -- SUVA PROTOCOL
% QuantumVault Smart Contract
% Pour : William Schulz
% Par  : Mohamed Lamine MBENGUE
% Date : Mars 2026
%================================================
\documentclass[11pt,a4paper]{article}

\usepackage[utf8]{inputenc}
\usepackage[T1]{fontenc}
\usepackage[french]{babel}
\usepackage{geometry}
\geometry{margin=2.5cm}

\usepackage{amsmath,amssymb,amsthm}
\usepackage{booktabs}
\usepackage{xcolor}
\usepackage{listings}
\usepackage{hyperref}
\usepackage{graphicx}
\usepackage{tikz}
\usetikzlibrary{arrows.meta, positioning, shapes.geometric, calc}
\usepackage{fancyhdr}
\usepackage{float}
\usepackage{array}
\usepackage{multirow}
\usepackage{tcolorbox}
\tcbuselibrary{skins, breakable}

%--- Couleurs ---
\definecolor{suva-blue}{RGB}{0, 82, 165}
\definecolor{suva-green}{RGB}{0, 150, 90}
\definecolor{suva-orange}{RGB}{220, 100, 0}
\definecolor{code-bg}{RGB}{245, 245, 245}
\definecolor{pass-green}{RGB}{0, 140, 70}
\definecolor{fail-red}{RGB}{180, 0, 0}

%--- En-tete/Pied de page ---
\pagestyle{fancy}
\fancyhf{}
\rhead{\textcolor{suva-blue}{\textbf{SuVa Protocol --- Rapport Phase 3}}}
\lhead{\textit{Confidentiel}}
\rfoot{\thepage}
\lfoot{\textit{Mars 2026}}

%--- Styles de code ---
\lstdefinestyle{solidity}{
  language=C,
  basicstyle=\ttfamily\small,
  backgroundcolor=\color{code-bg},
  keywordstyle=\color{suva-blue}\bfseries,
  commentstyle=\color{gray}\itshape,
  numbers=left, numberstyle=\tiny\color{gray},
  frame=single, breaklines=true, captionpos=b
}
\lstdefinestyle{python}{
  language=Python,
  basicstyle=\ttfamily\small,
  backgroundcolor=\color{code-bg},
  keywordstyle=\color{suva-blue}\bfseries,
  commentstyle=\color{suva-green}\itshape,
  numbers=left, numberstyle=\tiny\color{gray},
  frame=single, breaklines=true, captionpos=b
}
\lstdefinestyle{bash}{
  language=bash,
  basicstyle=\ttfamily\small,
  backgroundcolor=\color{code-bg},
  keywordstyle=\color{suva-orange}\bfseries,
  commentstyle=\color{gray}\itshape,
  numbers=none,
  frame=single, breaklines=true, captionpos=b
}

%--- Boites personnalisees ---
\newtcolorbox{resultbox}[1][]{
  enhanced,colback=suva-green!8,colframe=suva-green!60,
  fonttitle=\bfseries,title=#1,arc=3mm
}
\newtcolorbox{warningbox}[1][]{
  enhanced,colback=suva-orange!8,colframe=suva-orange!60,
  fonttitle=\bfseries,title=#1,arc=3mm
}
\newtcolorbox{infobox}[1][]{
  enhanced,colback=suva-blue!6,colframe=suva-blue!50,
  fonttitle=\bfseries,title=#1,arc=3mm
}

%================================================
\begin{document}
%================================================

\begin{titlepage}
\centering
\vspace*{2cm}
{\Huge\bfseries\textcolor{suva-blue}{SuVa Protocol}\par}
\vspace{0.5cm}
{\Large\textit{Stablecoin Quantique sur EVM}\par}
\vspace{1.5cm}
\rule{\linewidth}{0.5pt}
\vspace{0.5cm}
{\LARGE\bfseries Rapport Technique --- Phase 3\par}
\vspace{0.3cm}
{\large QuantumVault.sol --- Coffre-fort EVM garde par signature post-quantique Falcon\par}
\vspace{0.5cm}
\rule{\linewidth}{0.5pt}
\vspace{2cm}

\begin{tabular}{ll}
\textbf{Client :}    & William Schulz \\[4pt]
\textbf{Auteur :}    & Mohamed Lamine MBENGUE \\[4pt]
\textbf{Portfolio :} & \href{https://portfolio-alamine.web.app}{\texttt{portfolio-alamine.web.app}} \\[4pt]
\textbf{Date :}      & Mars 2026 \\[4pt]
\textbf{Version :}   & 1.0 --- Final \\[4pt]
\textbf{Statut :}    & \textcolor{pass-green}{\textbf{COMPLET}} \\[4pt]
\textbf{Codebase :}  & ETHDILITHIUM-main (fork ZKNox) \\
\end{tabular}

\vfill
\begin{tcolorbox}[enhanced,colback=suva-blue!8,colframe=suva-blue,arc=4mm]
\centering\textit{Document confidentiel. Ce rapport couvre les livrables de la Phase 3 :\\
contrat QuantumVault.sol, token ERC-20 qUSDT, enregistrement de cle Falcon on-chain,\\
tests d'integration bout-en-bout, et benchmarks de gas.}
\end{tcolorbox}
\end{titlepage}

\tableofcontents
\newpage

%================================================
\section{Resume Executif}
%================================================

\begin{resultbox}[Phase 3 --- Tous les objectifs atteints]
\begin{itemize}
  \item[$\checkmark$] \textbf{QuantumVault.sol} deploye et teste : custody USDT protegee par signature Falcon
  \item[$\checkmark$] \textbf{qUSDT.sol} ERC-20 : minte au depot (1:1), brule au retrait
  \item[$\checkmark$] \textbf{Enregistrement de cle Falcon on-chain} : hash de la cle publique stocke par adresse
  \item[$\checkmark$] \textbf{Protection anti-rejeu} : nonce par utilisateur incremente a chaque retrait
  \item[$\checkmark$] \textbf{14 tests de logique vault} passant avec MockFalcon
  \item[$\checkmark$] \textbf{1 test bout-en-bout reel} passant avec le vrai verificateur ZKNOX\_ethfalcon
  \item[$\checkmark$] \textbf{26 forge tests au total}, 0 echec, toutes phases confondues
  \item[$\checkmark$] Benchmarks de gas mesures : overhead vault seulement 174\,133 gas
\end{itemize}
\end{resultbox}

\vspace{0.5cm}

La Phase 3 livre le coeur du protocole SuVa : un smart contract qui detient des USDT et exige
une signature Falcon-256 valide pour autoriser tout retrait. Aucune ECDSA. Aucun admin de
confiance. Seul le detenteur de la cle privee Falcon correspondante peut deplacer les fonds.

La logique du vault (depot, enregistrement de cle, nonce, controle de retrait) n'ajoute que
\textbf{174\,133 gas} d'overhead. Le cout dominant reste le verificateur Falcon a \textbf{46,7M gas}
(multiplication scolaire sans NTT). La Phase 4 introduira l'optimisation NTT avec un objectif
de \textbf{$\sim$4M gas}.

%================================================
\section{Architecture}
%================================================

\subsection{Vue d'ensemble}

\begin{center}
\begin{tikzpicture}[
  node distance=1.8cm,
  box/.style={rectangle,draw,rounded corners,minimum width=3.2cm,minimum height=1cm,
              align=center,font=\small},
  arrow/.style={-Stealth,thick}
]
  \node[box, fill=suva-blue!15]   (user)   at (0,0)   {Utilisateur\\(Python)};
  \node[box, fill=suva-green!15]  (vault)  at (5.5,0) {QuantumVault.sol};
  \node[box, fill=suva-orange!15, above right=0.7cm and 2.5cm of vault] (falcon) {ZKNOX\_ethfalcon};
  \node[box, fill=suva-blue!10,   below right=0.7cm and 2.5cm of vault] (qusdt)  {qUSDT.sol (ERC-20)};
  \node[box, fill=gray!15, below=1.5cm of vault] (usdt) {MockERC20 (USDT)};

  \draw[arrow] (user)  -- node[above,font=\footnotesize]{depot / retrait} (vault);
  \draw[arrow] (vault) -- node[above,font=\footnotesize]{verify(h,sig,msg)} (falcon);
  \draw[arrow] (vault) -- node[above,font=\footnotesize]{mint / burn} (qusdt);
  \draw[arrow] (vault) -- node[left, font=\footnotesize]{transferFrom / transfer} (usdt);
\end{tikzpicture}
\end{center}

\begin{center}
\begin{tabular}{lll}
\toprule
\textbf{Contrat} & \textbf{Role} & \textbf{Etat cle} \\
\midrule
\texttt{QuantumVault.sol}  & Detient USDT, controle retraits  & \texttt{pkHash}, \texttt{deposits}, \texttt{nonces} \\
\texttt{qUSDT.sol}         & Token ERC-20 de recu             & Balances (ERC-20 herite) \\
\texttt{ZKNOX\_ethfalcon}  & Verifie signatures Falcon-256    & Sans etat \\
\texttt{MockERC20}         & Simulation USDT pour les tests   & Balances \\
\bottomrule
\end{tabular}
\end{center}

%================================================
\section{Conception des Smart Contracts}
%================================================

\subsection{QuantumVault.sol}

\subsubsection{Pattern Interface pour le Verificateur}

\begin{lstlisting}[style=solidity, caption={Interface IFalconVerifier}]
interface IFalconVerifier {
    function verify(
        uint256[256] calldata h,
        bytes calldata salt,
        bytes calldata s1,
        bytes calldata message
    ) external view returns (bool);
}
\end{lstlisting}

\begin{infobox}[Decision de conception : Pattern Interface]
Ce pattern permet d'injecter un \texttt{MockFalcon} dans les tests (qui retourne toujours
\texttt{true} ou \texttt{false}). Cela separe le test de la logique vault du test de la
verification cryptographique. Un futur verificateur NTT pourra etre branche sans modifier
le vault.
\end{infobox}

\subsubsection{Variables d'etat}

\begin{lstlisting}[style=solidity, caption={Stockage QuantumVault}]
IERC20           public immutable usdt;
IqUSDT           public immutable qusdt;
IFalconVerifier  public immutable falcon;

// keccak256(abi.encodePacked(h)) par utilisateur
mapping(address => bytes32) public pkHash;

// Montant USDT depose par utilisateur
mapping(address => uint256) public deposits;

// Nonce anti-rejeu par utilisateur
mapping(address => uint256) public nonces;
\end{lstlisting}

\subsubsection{Stockage de la Cle : Hash vs Cle Complete}

\begin{infobox}[Decision : Stocker le hash de la cle publique]
La cle publique Falcon-256 \texttt{h} est un tableau de 256 \texttt{uint256} = \textbf{8\,192 octets}.\\[4pt]
\textbf{Option A --- Stocker la cle complete :} 256 SSTORE $\approx$ \textbf{5,1M gas}\\
\textbf{Option B --- Stocker \texttt{keccak256(h)} :} 1 SSTORE $\approx$ \textbf{20\,000 gas}\\[4pt]
On choisit l'Option B. La cle complete est passee en \texttt{calldata} a chaque retrait.
Le vault verifie \texttt{keccak256(abi.encodePacked(h)) == pkHash[user]} avant de verifier la
signature. C'est cryptographiquement equivalent.
\end{infobox}

\subsubsection{Format du Message Signe}

\begin{lstlisting}[style=solidity, caption={withdrawMessage() --- le payload signe}]
function withdrawMessage(
    address user,
    uint256 amount,
    address to
) public view returns (bytes memory) {
    return abi.encodePacked(
        "SuVa:withdraw:",   // separateur de domaine
        block.chainid,      // ID de la chaine
        address(this),      // adresse du vault
        nonces[user],       // nonce anti-rejeu
        amount,             // montant exact
        to                  // adresse destinataire
    );
}
\end{lstlisting}

\begin{center}
\begin{tabular}{lll}
\toprule
\textbf{Champ} & \textbf{Valeur} & \textbf{Attaque prevenue} \\
\midrule
\texttt{"SuVa:withdraw:"} & Separateur domaine    & Rejeu inter-applications \\
\texttt{block.chainid}    & ex. 1 (mainnet)        & Rejeu inter-chaines \\
\texttt{address(this)}    & Adresse du vault       & Rejeu inter-contrats \\
\texttt{nonces[user]}     & Incremente par retrait & Rejeu de transaction \\
\texttt{amount}           & Montant USDT exact     & Manipulation montant \\
\texttt{to}               & Adresse destinataire   & Manipulation destination \\
\bottomrule
\end{tabular}
\end{center}

\subsubsection{Fonctions Principales}

\begin{lstlisting}[style=solidity, caption={registerKey}]
function registerKey(uint256[256] calldata h) external {
    pkHash[msg.sender] = keccak256(abi.encodePacked(h));
    emit KeyRegistered(msg.sender, pkHash[msg.sender]);
}
\end{lstlisting}

\begin{lstlisting}[style=solidity, caption={deposit}]
function deposit(uint256 amount) external {
    require(pkHash[msg.sender] != bytes32(0), "Aucune cle enregistree");
    require(amount > 0, "Montant nul");
    usdt.transferFrom(msg.sender, address(this), amount);
    deposits[msg.sender] += amount;
    qusdt.mint(msg.sender, amount);
    emit Deposited(msg.sender, amount);
}
\end{lstlisting}

\begin{lstlisting}[style=solidity, caption={withdraw (Verifications-Effets-Interactions)}]
function withdraw(
    uint256[256] calldata h,
    bytes calldata salt,
    bytes calldata s1,
    uint256 amount,
    address to
) external {
    // VERIFICATIONS
    require(amount > 0, "Montant nul");
    require(deposits[msg.sender] >= amount, "Solde insuffisant");
    require(
        keccak256(abi.encodePacked(h)) == pkHash[msg.sender],
        "Mauvaise cle publique"
    );
    bytes memory message = withdrawMessage(msg.sender, amount, to);
    require(
        falcon.verify(h, salt, s1, message),
        "Signature Falcon invalide"
    );
    // EFFETS
    deposits[msg.sender] -= amount;
    nonces[msg.sender]++;
    // INTERACTIONS
    qusdt.burn(msg.sender, amount);
    usdt.transfer(to, amount);
    emit Withdrawn(msg.sender, to, amount);
}
\end{lstlisting}

\begin{infobox}[Securite : Pattern Verifications-Effets-Interactions (CEI)]
Toutes les verifications et mises a jour d'etat se font \textit{avant} les appels externes.
Cela previent les attaques par reentrance.
\end{infobox}

\subsection{qUSDT.sol}

\begin{lstlisting}[style=solidity, caption={qUSDT --- token de recu quantique}]
contract qUSDT is ERC20 {
    address public vault;
    constructor() ERC20("Quantum USDT", "qUSDT") {}

    function setVault(address _vault) external {
        require(vault == address(0), "Deja configure");
        vault = _vault;
    }
    modifier onlyVault() {
        require(msg.sender == vault, "Pas le vault");
        _;
    }
    function mint(address to, uint256 amount) external onlyVault {
        _mint(to, amount);
    }
    function burn(address from, uint256 amount) external onlyVault {
        _burn(from, amount);
    }
}
\end{lstlisting}

%================================================
\section{Couverture des Tests}
%================================================

\subsection{Sortie forge test}

\begin{lstlisting}[style=bash, caption={Sortie complete forge test -vv}]
Ran 14 tests for test/QuantumVault.t.sol:QuantumVaultTest
[PASS] testDeposit()                     (gas: 402 008)
[PASS] testDepositRequiresKey()          (gas:  34 488)
[PASS] testDepositZeroReverts()          (gas: 257 388)
[PASS] testNonceIncrement()              (gas: 1 142 879)
[PASS] testQUSDTOnlyVaultBurn()          (gas: 395 122)
[PASS] testQUSDTOnlyVaultMint()          (gas:  34 832)
[PASS] testQUSDTSetVaultOnlyOnce()       (gas:  32 040)
[PASS] testRegisterKey()                 (gas: 261 475)
[PASS] testWithdraw()                    (gas: 801 715)
[PASS] testWithdrawInsufficientBalance() (gas: 756 598)
[PASS] testWithdrawInvalidSig()          (gas: 778 281)
[PASS] testWithdrawMessageEncoding()     (gas: 378 441)
[PASS] testWithdrawWrongKey()            (gas: 756 762)
[PASS] testWithdrawZeroReverts()         (gas: 756 410)
Suite result: ok. 14 passed; 0 failed

Ran 1 test for test/QuantumVaultReal.t.sol:QuantumVaultRealTest
[PASS] testRealFalconWithdraw()          (gas: 47 208 810)
Suite result: ok. 1 passed; 0 failed

Ran 10 test suites: 26 passed, 0 failed, 0 skipped
\end{lstlisting}

\begin{resultbox}[26/26 tests forge passants --- 0 echec]
Tous les tests Phase 1, Phase 2 et Phase 3 passent simultanement. Aucune regression.
\end{resultbox}

\subsection{Tableau des tests}

\begin{center}
\begin{tabular}{llc}
\toprule
\textbf{Test} & \textbf{Ce qu'il verifie} & \textbf{Resultat} \\
\midrule
\texttt{testRegisterKey}                 & pkHash stocke, evenement emis & \textcolor{pass-green}{\textbf{OK}} \\
\texttt{testDeposit}                     & USDT entrant, qUSDT minte 1:1 & \textcolor{pass-green}{\textbf{OK}} \\
\texttt{testDepositRequiresKey}          & Revert si aucune cle enregistree & \textcolor{pass-green}{\textbf{OK}} \\
\texttt{testDepositZeroReverts}          & Revert sur montant nul & \textcolor{pass-green}{\textbf{OK}} \\
\texttt{testWithdraw}                    & Retrait complet, USDT rendu & \textcolor{pass-green}{\textbf{OK}} \\
\texttt{testWithdrawInvalidSig}          & Fausse sig $\Rightarrow$ revert & \textcolor{pass-green}{\textbf{OK}} \\
\texttt{testWithdrawInsufficientBalance} & Retrait $>$ depot $\Rightarrow$ revert & \textcolor{pass-green}{\textbf{OK}} \\
\texttt{testNonceIncrement}              & Nonce s'incremente, ancienne sig rejete & \textcolor{pass-green}{\textbf{OK}} \\
\texttt{testWithdrawWrongKey}            & Mauvais \texttt{h} $\Rightarrow$ revert & \textcolor{pass-green}{\textbf{OK}} \\
\texttt{testWithdrawZeroReverts}         & Montant nul $\Rightarrow$ revert & \textcolor{pass-green}{\textbf{OK}} \\
\texttt{testWithdrawMessageEncoding}     & Message interne coherent & \textcolor{pass-green}{\textbf{OK}} \\
\texttt{testQUSDTSetVaultOnlyOnce}       & Second \texttt{setVault()} revert & \textcolor{pass-green}{\textbf{OK}} \\
\texttt{testQUSDTOnlyVaultMint}          & Mint direct non-vault $\Rightarrow$ revert & \textcolor{pass-green}{\textbf{OK}} \\
\texttt{testQUSDTOnlyVaultBurn}          & Burn direct non-vault $\Rightarrow$ revert & \textcolor{pass-green}{\textbf{OK}} \\
\midrule
\texttt{testRealFalconWithdraw}          & Bout-en-bout : vraie cle Falcon-256 & \textcolor{pass-green}{\textbf{OK}} \\
\bottomrule
\end{tabular}
\end{center}

\subsection{Structure du Test Bout-en-Bout}

Le test \texttt{testRealFalconWithdraw} connecte la pile complete :

\begin{lstlisting}[style=solidity, caption={QuantumVaultReal.t.sol (simplifie)}]
contract QuantumVaultRealTest is Test {
    // Vrais contrats
    ZKNOX_ethfalcon verifier;
    QuantumVault    vault;
    qUSDT           qusdtToken;
    MockERC20       usdtToken;

    // Vrais vecteurs de test Falcon-256 (generes par script Python)
    int16[] h;    // cle publique (256 coefficients)
    int16[] salt; // nonce de signature
    int16[] s1;   // polynome de signature

    function setUp() public {
        // Deploiement de la pile complete avec le vrai verificateur
        verifier  = new ZKNOX_ethfalcon();
        qusdtToken = new qUSDT();
        usdtToken  = new MockERC20();
        vault = new QuantumVault(address(verifier),
                                 address(usdtToken),
                                 address(qusdtToken));
        qusdtToken.setVault(address(vault));

        // Chargement des vecteurs de test Falcon-256 depuis Python
        _loadTestVectors();
    }

    function testRealFalconWithdraw() public {
        // 1. Enregistrement de la vraie cle Falcon
        vault.registerKey(h);

        // 2. Depot de 1 000 000 USDT
        usdtToken.mint(address(this), 1_000_000e6);
        usdtToken.approve(address(vault), 1_000_000e6);
        vault.deposit(1_000_000e6);

        // 3. Signature du message de retrait (generee par Python)
        // La sig couvre : chainId + vault + nonce=0 + montant + to
        vault.withdraw(h, salt, s1, 1_000_000e6, address(this));

        // 4. Verification : USDT rendu, qUSDT brule
        assertEq(usdtToken.balanceOf(address(this)), 1_000_000e6);
        assertEq(qusdtToken.balanceOf(address(this)), 0);
    }
}
\end{lstlisting}

\newpage
%================================================
\section{Analyse du Gas}
%================================================

\subsection{Benchmark complet}

\begin{center}
\begin{tabular}{llr}
\toprule
\textbf{Contrat} & \textbf{Fonction} & \textbf{Gas moyen} \\
\midrule
QuantumVault & \texttt{registerKey}                   &     117\,540 \\
QuantumVault & \texttt{deposit}                       &     112\,218 \\
QuantumVault & \texttt{withdraw} (MockFalcon)         &     210\,410 \\
QuantumVault & \texttt{withdraw} (ETHFalcon reel)     &  46\,839\,799 \\
\midrule
ZKNOX\_ethfalcon    & \texttt{verify} (Falcon-256 scolaire) & 46\,665\,666 \\
ZKNOX\_ethdilithium & \texttt{verify} (ETHDilithium)        &  6\,481\,497 \\
ZKNOX\_dilithium    & \texttt{verify} (Dilithium standard)  & 13\,370\,438 \\
\midrule
qUSDT     & \texttt{mint}  & 24\,084 \\
qUSDT     & \texttt{burn}  & 24\,039 \\
MockERC20 & \texttt{mint}  & 65\,878 \\
\bottomrule
\end{tabular}
\end{center}

\subsection{Decomposition de l'overhead}

\begin{lstlisting}[style=bash, caption={Decomposition de l'overhead du vault}]
Cout total retrait  =  Verif. Falcon  +  Logique vault
                    =   46 665 666   +    174 133
                    =   46 839 799 gas

Logique vault seule (MockFalcon) = 210 410 gas  <- overhead SuVa
\end{lstlisting}

\begin{infobox}[Interpretation]
Le vault SuVa n'ajoute que \textbf{210K gas} d'overhead par retrait. Ce cout couvre :
2 SSTORE USDT/qUSDT, 1 SSTORE nonce, verification du hash (keccak + SLOAD), emission d'evenement,
et les couts calldata pour passer \texttt{h}.\\[4pt]
\textbf{L'overhead du vault est negligeable par rapport au cout de verification Falcon.}
Tout effort d'optimisation doit etre dirige vers le verificateur.
\end{infobox}

\subsection{Projection Phase 4 avec NTT}

\begin{lstlisting}[style=bash, caption={Projection de gas Phase 4}]
ETHFalcon actuel (scolaire) :    46 665 666 gas
ETHFalcon NTT Phase 4 (cible) :  ~4 000 000 gas
Overhead vault :                     210 410 gas
                                  -----------
Total retrait vault (Phase 4) :   ~4 210 410 gas

Points de reference :
  Transaction ECDSA standard :          21 000 gas
  Verification ETHDilithium :        6 481 497 gas
  ETHFalcon scolaire :               46 665 666 gas
  ETHFalcon NTT (cible Phase 4) :     4 000 000 gas
\end{lstlisting}

L'optimisation NTT exploite la structure de l'anneau NTRU de Falcon ($q = 12289$) pour remplacer
la multiplication polynomiale scolaire ($O(n^2)$) par une multiplication basee NTT ($O(n \log n)$),
visant une reduction d'environ $11\times$ du gas de verification.

\subsection{Cout sur L2}

\begin{lstlisting}[style=bash, caption={Projections de cout sur L2 (ETH = 4 000$)}]
--- Ethereum Mainnet a 30 gwei ---
ETHFalcon scolaire : 46,7M x 30 gwei = 1,401 ETH ~ 5 600$ / tx  [INACCEPTABLE]
ETHFalcon NTT Ph4 :   4,2M x 30 gwei = 0,126 ETH ~   504$ / tx  [COUTEUX]

--- Arbitrum a 0,1 gwei (300x moins cher que mainnet) ---
ETHFalcon NTT Ph4 :   4,2M x 0,1 gwei = 0,00042 ETH ~ 1,68$ / tx  [VIABLE]

--- Base a 0,01 gwei (3000x moins cher que mainnet) ---
ETHFalcon NTT Ph4 :   4,2M x 0,01 gwei = 0,000042 ETH ~ 0,17$ / tx  [VIABLE EN PRODUCTION]
\end{lstlisting}

\begin{infobox}[Recommandation de deploiement]
La Phase 3 etablit que SuVa est techniquement correct sur mainnet. Cependant, le deploiement
en production devrait cibler \textbf{Base} ou \textbf{Arbitrum} apres l'optimisation NTT Phase 4,
atteignant un cout inferieur a 1\$ par transaction. Comparable aux couts de signature hardware
wallet, acceptable pour un stablecoin quantique-resistant.
\end{infobox}

\newpage
%================================================
\section{Flux Bout-en-Bout}
%================================================

\subsection{Diagramme de Sequence}

\begin{figure}[H]
\centering
\begin{tikzpicture}[
  every node/.style={font=\small},
  actor/.style={rectangle, draw, rounded corners, fill=suva-blue!12, minimum width=2.4cm,
                minimum height=0.8cm, align=center},
  msg/.style={-Stealth, thick},
  ret/.style={-Stealth, thick, dashed, gray},
  lifeline/.style={dashed, gray, thick}
]

% Acteurs
\node[actor] (user)   at (0,0)   {Utilisateur\\(Python)};
\node[actor] (vault)  at (4.5,0) {QuantumVault};
\node[actor] (verif)  at (9,0)   {ZKNOX\_ethfalcon};
\node[actor] (qusdtN) at (13,0)  {qUSDT / USDT};

% Lignes de vie
\foreach \x in {0, 4.5, 9, 13} {
  \draw[lifeline] (\x, -0.4) -- (\x, -14.5);
}

% --- registerKey ---
\draw[msg] (0,-1.0) -- node[above, font=\footnotesize]{registerKey(h)} (4.5,-1.0);
\node[right, font=\footnotesize, text=suva-blue] at (4.6,-1.5) {pkHash[user] = keccak256(h)};

% --- deposit ---
\draw[msg] (0,-2.5) -- node[above, font=\footnotesize]{deposit(1 000 000)} (4.5,-2.5);
\draw[msg] (4.5,-3.0) -- node[above, font=\footnotesize]{\small usdt.transferFrom(user, vault, amt)} (13,-3.0);
\draw[msg] (4.5,-3.7) -- node[above, font=\footnotesize]{\small mint(user, amt)} (13,-3.7);
\node[right, font=\footnotesize, text=suva-green] at (13.1,-4.1) {qUSDT minte};

% --- signature locale ---
\node[left, font=\footnotesize\itshape, text=gray] at (-0.1,-5.2) {sign(msg) localement};
\draw[msg, bend left=30] (0,-5.0) to node[above, font=\footnotesize\itshape]{(hors-chaine)} (0,-5.6);

% --- withdraw ---
\draw[msg] (0,-6.5) -- node[above, font=\footnotesize]{withdraw(h, salt, s1, amt, to)} (4.5,-6.5);
\node[right, font=\footnotesize, text=suva-blue] at (4.6,-7.1) {verif pkHash[user] == keccak256(h)};
\node[right, font=\footnotesize, text=suva-blue] at (4.6,-7.6) {construit withdrawMessage(...)};
\draw[msg] (4.5,-8.2) -- node[above, font=\footnotesize]{verify(h, salt, s1, msg)} (9,-8.2);
\draw[ret] (9,-8.8) -- node[above, font=\footnotesize\itshape]{retourne true} (4.5,-8.8);
\node[right, font=\footnotesize, text=suva-blue] at (4.6,-9.4) {nonces[user]++};
\draw[msg] (4.5,-10.0) -- node[above, font=\footnotesize]{burn(user, amt)} (13,-10.0);
\draw[msg] (4.5,-10.7) -- node[above, font=\footnotesize]{usdt.transfer(to, amt)} (13,-10.7);
\node[right, font=\footnotesize, text=suva-green] at (13.1,-11.2) {USDT rendu};

\end{tikzpicture}
\caption{Flux SuVa bout-en-bout : enregistrement cle, depot, et retrait gate par Falcon.}
\end{figure}

\subsection{Script Python de Signature}

Le test bout-en-bout s'appuie sur un script Python qui genere un keypair Falcon-256 valide
et signe le message de retrait :

\begin{lstlisting}[style=python, caption={Signature du message de retrait (Python)}]
from falcon import SecretKey, PublicKey
import eth_abi

# 1. Generation de cle
sk = SecretKey(256)
pk = PublicKey(sk)

# 2. Construction du message de retrait (doit correspondre a Solidity withdrawMessage())
chain_id   = 31337      # Anvil / Foundry chaine locale
vault_addr = "0x..."    # adresse du vault deploye
nonce      = 0
amount     = 1_000_000_000_000  # 1M USDT (6 decimales)
to         = "0x..."    # destinataire

msg = b"SuVa:withdraw:" + \
      chain_id.to_bytes(32, 'big') + \
      bytes.fromhex(vault_addr[2:]) + \
      nonce.to_bytes(32, 'big') + \
      amount.to_bytes(32, 'big') + \
      bytes.fromhex(to[2:])

# 3. Signature
sig = sk.sign(msg)
salt, s1 = sig.nonce, sig.s1  # tableaux int16[]

# 4. Export comme vecteur de test Solidity
print(f"h    = {list(pk.h)}")
print(f"salt = {list(salt)}")
print(f"s1   = {list(s1)}")
\end{lstlisting}

\newpage
%================================================
\section{Decisions de Conception et Compromis}
%================================================

Cette section repond en detail aux questions de conception soulevees dans le brief de William Schulz.

\subsection{Cle publique complete on-chain vs. hash ?}

\textbf{Decision : stocker \texttt{keccak256(h)}, passer la cle complete \texttt{h} en calldata.}

\begin{center}
\begin{tabular}{lrr}
\toprule
\textbf{Option} & \textbf{Gas (registerKey)} & \textbf{Gas (calldata retrait)} \\
\midrule
Stocker la cle complete (8\,192 octets) & $\sim$5\,120\,000 & $\sim$0 \\
Stocker le hash, passer en calldata & $\sim$20\,000 & $\sim$131\,000 \\
\midrule
\textbf{Economies (enregistrement)} & \textbf{5\,100\,000 gas} & --- \\
\bottomrule
\end{tabular}
\end{center}

Argument de securite : \texttt{keccak256(h)} est un engagement par resistance aux collisions sur \texttt{h}.
Un utilisateur qui enregistre \texttt{keccak256(h)} est lie a cette cle specifique. Au retrait,
il passe la cle complete en calldata ; le vault recalcule le hash et rejette toute substitution.
La resistance a la preimage de Keccak-256 garantit qu'aucun attaquant ne peut trouver un
\texttt{h'} tel que \texttt{keccak256(h') == keccak256(h)}.

\subsection{Rotation de cle ?}

\textbf{Comportement Phase 3 :} tout utilisateur peut appeler \texttt{registerKey()} a tout moment pour
ecraser son \texttt{pkHash}. Aucune signature Falcon n'est requise --- seule la possession de la
cle privee Ethereum est necessaire.

\textbf{Limitation connue :} un attaquant qui compromet la cle ECDSA de l'utilisateur (mais pas sa
cle Falcon) peut ecraser le \texttt{pkHash} avec une cle qu'il controle, puis attendre que
l'utilisateur depose. C'est le probleme de couplage cle ECDSA-Falcon.

\textbf{Mitigation Phase 4 :} la rotation de cle exigera une signature Falcon valide de la cle
actuellement enregistree, protegant la rotation avec une authentification quantique-resistante.

\subsection{Perte de la cle privee Falcon}

\textbf{Aucun mecanisme de recuperation n'existe.} Si un utilisateur perd sa cle privee Falcon-256 :

\begin{warningbox}[Aucun mecanisme de recuperation --- Par conception]
\begin{itemize}
  \item Il ne peut pas produire une signature de retrait valide.
  \item Aucun admin, aucun timelock, aucune recuperation sociale ne peut recuperer les fonds.
  \item Le solde USDT et le qUSDT de l'utilisateur restent definitivement irrecuperables.
\end{itemize}
C'est une decision de conception explicite et intentionnelle. Tout mecanisme de recuperation
cree un acteur privilegie qui pourrait etre compromis, contraint ou en collusion avec un
adversaire quantique. L'ensemble du modele de securite repose sur le fait que la cle privee
Falcon soit le \textit{seul} chemin d'autorisation.\\[4pt]
\textbf{Responsabilite utilisateur :} Sauvegardez votre cle privee Falcon.
\end{warningbox}

\subsection{Contrat evolutif ?}

\textbf{Phase 3 : pas de pattern proxy.} Le vault est immuable. Avantages :
\begin{itemize}
  \item Pas de cle admin avec pouvoir de mise a jour (reduit les hypotheses de confiance)
  \item Pas de risques de collision dans le stockage proxy (EIP-1967)
  \item Surface d'audit plus simple
  \item En cas de bug critique, un nouveau vault est deploye avec une notice aux utilisateurs
\end{itemize}

La Phase 4 pourra introduire un proxy evolutif, mais seulement si le mecanisme de mise a
jour est lui-meme gate par des signatures Falcon (les operations admin requi\`erent une
autorisation quantique-resistante).

%================================================
\section{Checklist des Livrables}
%================================================

\begin{center}
\begin{tabular}{clc}
\toprule
& \textbf{Livrable} & \textbf{Statut} \\
\midrule
$\checkmark$ & \texttt{QuantumVault.sol} implemente et teste & \textcolor{pass-green}{\textbf{FAIT}} \\
$\checkmark$ & \texttt{qUSDT.sol} ERC-20 (mint/burn) & \textcolor{pass-green}{\textbf{FAIT}} \\
$\checkmark$ & Enregistrement de cle Falcon on-chain & \textcolor{pass-green}{\textbf{FAIT}} \\
$\checkmark$ & Test d'integration bout-en-bout & \textcolor{pass-green}{\textbf{FAIT}} \\
$\checkmark$ & 14 tests logique vault (MockFalcon) & \textcolor{pass-green}{\textbf{FAIT}} \\
$\checkmark$ & 1 test Falcon reel bout-en-bout & \textcolor{pass-green}{\textbf{FAIT}} \\
$\checkmark$ & Benchmarks de gas mesures et documentes & \textcolor{pass-green}{\textbf{FAIT}} \\
$\checkmark$ & 26/26 tests totaux passants & \textcolor{pass-green}{\textbf{FAIT}} \\
$\checkmark$ & Deploiement sur Sepolia testnet & \textcolor{pass-green}{\textbf{FAIT}} \\
\bottomrule
\end{tabular}
\end{center}

\vspace{0.4cm}

\begin{resultbox}[Phase 3 : 9/9 objectifs completes --- LIVRE INTEGRALEMENT]
Tous les objectifs smart contract de la Phase 3 sont livres, testes et deployes sur le testnet Sepolia.
\end{resultbox}

\newpage
\subsection{Deploiement sur Sepolia Testnet}

Tous les contrats SuVa Phase 3 sont deployes et verifiables sur le testnet Ethereum Sepolia (Chain ID : 11155111).

\begin{center}
\begin{tabular}{llr}
\toprule
\textbf{Contrat} & \textbf{Adresse Sepolia} & \textbf{Bloc} \\
\midrule
\texttt{ZKNOX\_ethfalcon} & \texttt{0xb82B557177e96425046e1a8036F98EB02d07Cdeb} & 10554061 \\
\texttt{MockERC20 (USDT)} & \texttt{0x683e75234d1FD52008d6FBC7B2c671957C02D41a} & 10554061 \\
\texttt{qUSDT}            & \texttt{0x6E291d7c9121DA76697491f1211eF6fbb71df646} & 10554061 \\
\texttt{QuantumVault}     & \texttt{0x83773c48f72176982E90C0742361a52dAaDA1e19} & 10554061 \\
\bottomrule
\end{tabular}
\end{center}

\begin{center}
\begin{tabular}{ll}
\toprule
\textbf{Parametre} & \textbf{Valeur} \\
\midrule
Adresse deploieur   & \texttt{0xD7390D91139Bf5868DDdF3b6677708b30b78D210} \\
Tx de deploiement   & \texttt{0x1454d039072e5fc717c534ae5596cec5b191549c55c0ab87a38e17cfc0680ddc} \\
Gas total utilise   & 2\,796\,227 gas \\
Cout total          & 0,000003635 ETH \\
Reseau              & Ethereum Sepolia (Chain ID 11155111) \\
\bottomrule
\end{tabular}
\end{center}

\begin{infobox}[Verification sur Sepolia Etherscan]
Tous les contrats sont publiquement verifiables sur \texttt{sepolia.etherscan.io}.\\
Exemple : chercher \texttt{0x83773c48f72176982E90C0742361a52dAaDA1e19} pour inspecter QuantumVault.
\end{infobox}

%================================================
\section{Apercu Phase 4}
%================================================

\subsection{Prochaines Etapes}

\begin{center}
\begin{tabular}{clll}
\toprule
\textbf{Priorite} & \textbf{Tache} & \textbf{Metrique cible} & \textbf{Dependance} \\
\midrule
1 & ETHFalcon optimise NTT & $\sim$4M gas verify & Impl. NTT ZKNox \\
3 & \texttt{SuVaWallet.sol} (ERC-4337) & Double-sig : ECDSA + Falcon & Vault Phase 3 stable \\
4 & Rotation de cle gate par Falcon & \texttt{registerKey} requiert ancienne sig & Verificateur NTT Phase 4 \\
\bottomrule
\end{tabular}
\end{center}

\subsection{Raison de l'Optimisation NTT}

Le \texttt{ZKNOX\_ethfalcon} actuel utilise la multiplication polynomiale scolaire pour
l'equation de verification NTRU :
\[
  h \cdot s_1 + s_2 \equiv 0 \pmod{q}
\]
ou $h, s_1, s_2 \in \mathbb{Z}_q[X]/(X^n + 1)$, $q = 12289$, $n = 512$.

La multiplication scolaire coute $O(n^2) = O(512^2) = 262\,144$ multiplications de coefficients.
Puisque $q = 12289$ est un premier compatible NTT ($q = 12 \cdot 1024 + 1$), la multiplication
polynomiale peut etre remplacee par :
\[
  \text{NTT}^{-1}\bigl(\text{NTT}(h) \odot \text{NTT}(s_1)\bigr) + s_2
\]
a un cout de $O(n \log n) = O(512 \cdot 9) \approx 4\,608$ operations ---
une reduction theorique de $57\times$. L'overhead EVM pratique devrait ramener cela a
une reduction de gas d'environ $11\times$ (de 46,7M a $\sim$4M gas).

\subsection{Architecture ERC-4337 SuVaWallet}

Le wallet Phase 4 implementera l'interface ERC-4337 \texttt{IAccount} avec validation
a double signature :

\begin{lstlisting}[style=solidity, caption={Interface SuVaWallet.sol (apercu Phase 4)}]
contract SuVaWallet is IAccount {
    function validateUserOp(
        UserOperation calldata userOp,
        bytes32 userOpHash,
        uint256 missingAccountFunds
    ) external returns (uint256 validationData) {
        // 1. Verifier la signature ECDSA (proprietaire)
        // 2. Verifier la signature Falcon-256 (quantique-resistant)
        // Les deux doivent passer pour que l'operation se poursuive
    }
}
\end{lstlisting}

Cette approche a double signature offre une defense en profondeur : meme si les ordinateurs
quantiques cassent ECDSA a l'avenir, la signature Falcon reste valide. Et si un bug est
trouve dans le verificateur Falcon, la signature ECDSA sert de repli.

%================================================
\section{Conclusion}
%================================================

La Phase 3 livre la couche fonctionnelle centrale du protocole SuVa. Le contrat
\textbf{QuantumVault.sol} fournit :

\begin{enumerate}
  \item \textbf{Autorisation quantique-resistante :} les retraits exigent une signature Falcon-256
        valide, inviolable par les ordinateurs quantiques.
  \item \textbf{Logique vault efficace :} seulement 210K gas d'overhead, independant du cout
        du schema de signature.
  \item \textbf{Separation nette des responsabilites :} l'interface \texttt{IFalconVerifier}
        decouple la logique vault de l'implementation cryptographique, permettant une
        mise a niveau sans redeploiement du vault lors du passage au verificateur NTT Phase 4.
  \item \textbf{Couverture de tests complete :} 14 tests unitaires couvrant tous les chemins
        de revert, plus 1 test d'integration bout-en-bout avec la vraie cryptographie Falcon-256.
  \item \textbf{Profil de gas honnete :} le cout actuel de 46,7M gas est entierement compris
        et attribue au verificateur Falcon scolaire. L'optimisation NTT Phase 4 devrait
        le reduire a $\sim$4M gas, rendant le deploiement en production sur reseaux L2
        viable a moins de 1\$ par transaction.
\end{enumerate}

Le protocole SuVa est sur la bonne voie. Tous les objectifs smart contract de la Phase 3 sont
completes. Les fondations cryptographiques validees en Phase 1 (Dilithium), les outils
d'integration de la Phase 2 (keygen Falcon, validation croisee), et les contrats vault de la
Phase 3 forment un systeme coherent, teste et auditable pret pour la prochaine phase.

%================================================
\appendix
\section{Annexe A --- Environnement}
%================================================

\begin{center}
\begin{tabular}{ll}
\toprule
\textbf{Outil} & \textbf{Version} \\
\midrule
Python & 3.14 (environnement virtuel) \\
Foundry (forge) & Nightly (win32\_amd64) \\
Solidity & 0.8.25 \\
Version EVM & Prague \\
Falcon (Python) & Fork ZKNox (base NTRU) \\
ZKNOX\_ethfalcon & scolaire (Phase 3) \\
ZKNOX\_ethdilithium & Variante XOF Keccak-256 \\
OpenZeppelin & v4.x (base ERC-20) \\
\bottomrule
\end{tabular}
\end{center}

%================================================
\section{Annexe B --- Parametres Falcon-256}
%================================================

\begin{center}
\begin{tabular}{lrl}
\toprule
\textbf{Parametre} & \textbf{Valeur} & \textbf{Description} \\
\midrule
$n$ & 512 & Dimension de l'anneau ($X^{512} + 1$) \\
$q$ & 12289 & Module (compatible NTT : $q = 12 \cdot 1024 + 1$) \\
$\sigma$ & $\approx 165,74$ & Ecart-type gaussien \\
Taille pk & 897 octets & Format compact NIST \\
Taille sig & 666 octets & Format compact NIST \\
Taille tableau \texttt{h} & 256 int16 & Representation ZKNox \\
Niveau de securite & NIST I & $\geq 128$ bits classique, $\geq 64$ bits quantique \\
\bottomrule
\end{tabular}
\end{center}

%================================================
\section{Annexe C --- Structure du Depot (Phase 3)}
%================================================

\begin{center}
\begin{tabular}{ll}
\toprule
\textbf{Chemin} & \textbf{Contenu} \\
\midrule
\texttt{src/QuantumVault.sol} & Contrat vault principal \\
\texttt{src/qUSDT.sol} & Token de recu ERC-20 \\
\texttt{src/IFalconVerifier.sol} & Interface pour les backends de verification \\
\texttt{src/MockFalcon.sol} & Double de test du verificateur \\
\texttt{test/QuantumVault.t.sol} & 14 tests unitaires (MockFalcon) \\
\texttt{test/QuantumVaultReal.t.sol} & 1 test bout-en-bout (ZKNOX\_ethfalcon) \\
\texttt{python-ref/vault\_sign.py} & Script de signature Falcon-256 pour les tests \\
\texttt{rapports/rapport\_phase3\_EN.tex} & Version anglaise \\
\texttt{rapports/rapport\_phase3\_FR.tex} & Ce document \\
\bottomrule
\end{tabular}
\end{center}

%================================================
\end{document}
%================================================
